Per pair or trio: printed pages, pencils and one strip set from \textbf{week-51-materials.pdf} on card stock at actual size (two sheets for a table of two pairs). \textbf{One unit is 10 mm}; each main 0-16 ruler and each sheet row is 160 mm long. A set has two gray strips of each whole length from 1 to 15 units and two white add-ons each of 4 and 1 units for the shared-A page; every extreme in the keys is a whole length. To take a length from a band, lay a strip along it from 0 and check that the end lands inside; join strips end to end on a 0-16 ruler. Repeated A means two equal strips, or one strip used twice. Cuisenaire rods also fit. Cutting takes about 10 minutes per set, less with a paper trimmer; the sets are reusable. The kit is physically unrehearsed.

\textbf{Readiness.} K-1 needs matching endpoints and counting small lengths, with adult reading. Older children need addition/subtraction and comparison of intervals. Negative numbers are unnecessary: in every printed gap task, the designated upper/longer join stays to the right. Keep an adult at each of the three tables.

\textbf{Launch, 4 minutes.} Let children find the strips whose ends land inside a band. Show A from 2 to 3 and B from 1 to 2 using the printed visual: A=2 and B=2, joined gap-free from 0, end at 4, inside the A + B band. Ask for a shorter join and a longer one; the shortest is 3 and the longest 5. Ask, ``Can you choose allowed lengths that break my promise?''

\textbf{Flexible hour.} 0-5 arrival; 5-9 launch; 9-47 tables: P1-2 with partner tests, then P3 (older bands also P4) or return pages; 47-53 each pair shows one guarantee and the two settings at its ends; 53-60 reset. Keep uncertain ranges visibly distinct from one actual trial.
